\documentclass[10pt,a4paper]{amsart}
\usepackage[margin=19mm]{geometry}
\usepackage{amsmath,amssymb,mathtools}
\usepackage[scr=rsfs,cal=euler]{mathalfa}
\usepackage{xfrac}
\usepackage[unicode,hidelinks,bookmarks=false]{hyperref}
\newcommand{\ie}{i.e.\ }
\newcommand{\cf}{cf.\ }
\newcommand{\resp}{respectively\ }
\newcommand{\Subsection}{\subsection}
\providecommand{\nobreakdash}{\nobreak\hbox{-}}
\setlength{\emergencystretch}{2em}
\multlinegap0pt
\usepackage{amssymb}
\newtheorem{theorem}{Theorem}
\title{To cover a permutohedron}
\author{Bochao Kong, Ji Zeng}
\date{Source: arXiv:2509.13877v3 (CC BY 4.0)}
\begin{document}
\maketitle

\noindent\emph{Source and licence.} To cover a permutohedron. Version arXiv:2509.13877v3 (CC BY 4.0), distributed by arXiv under the Creative Commons Attribution 4.0 International licence, http://creativecommons.org/licenses/by/4.0/, as recorded in the arXiv licence metadata for this exact version and shown on the abstract page https://arxiv.org/abs/2509.13877v3. Full licence notice: \texttt{licenses/arxiv-2509.13877-CC-BY-4.0.txt}.\par\smallskip

\noindent\textit{Editorial scope.} Counted unit: the article's theorem thm:main (source0.tex lines 34--36). The article's complete original proof of it is delivered with it (source0.tex lines 60--69, one \texttt{proof} environment of 10 lines, which the article places away from the statement). No mathematics is written, repaired or re-derived: the statement and the proof are the article's own lines, in the article's own notation, and no retained line is substituted or re-flowed.  Retained uncounted as the source-local context the proof needs: lines 54--56.  Nothing was omitted from any retained span, so the package carries no omission record.  No retained line contains a citation, so the wrapper carries no bibliography and no bibliography entry.  No retained line contains a figure environment, an image-inclusion command or drawing code, so no image is delivered and the package declares no asset dependency.  Editorial formatting only, in addition: an AMS \texttt{amsart} preamble that restates the article's own declarations and macros used by the retained lines, namely the environments \texttt{theorem} on separate counters, together with the standard packages those lines need.  The retained environments print in this document as Theorem 1 in order of appearance, whereas the article prints the same items with the numbering of its own full document.  The complete native source of the article as distributed in the arXiv e-print for this exact version is bundled unchanged as \texttt{sources/185194/source0.tex}.
\begin{theorem}\label{general}
  Assume that $F_A$ has $n-1$ distinct critical values. Then every affine hyperplane $H\subset \mathbb{R}^n$ with $H\neq H_A$ contains at most $(n-1)!$ vertices of $P_A$. In particular, any vertex cover of $P_A$ has size at least $n$.
\end{theorem}

\begin{theorem}\label{thm:main}
  If $n\ge 3$ is odd, then every affine hyperplane $H\subset\mathbb{R}^n$ with $H\neq H_n$ contains at most $(n-1)!$ vertices of $P_n$.
\end{theorem}

\begin{proof}[Proof of Theorem~\ref{thm:main}]

  Let $n \geq 3$ be odd and write $F_{n} = F_{[n]}$. According to Theorem~\ref{general}, it suffices to verify that $F_n$ has $n-1$ distinct critical values. We consider the real polynomial $G_n(x) = \prod_{j\in [n]} (x-j)$ and notice that $G_n(x) = F_n(x) - n!$. Hence, the critical points of $G_n$ coincide with those of $F_n$, and it suffices to prove $G_n$ has $n-1$ distinct critical values.

  Now, let $v_j$ be the extreme value of $G_n$ on the interval $(j,j+1)$ for $j\in [n-1]$. Notice that $v_1,\dots,v_{n-1}$ are the critical values of $G_n$. Since $n$ is odd, the graph of the function $G_n$ is symmetric with respect to the point $\frac{n+1}{2}$ on the $x$-axis. Therefore, it suffices to prove $|v_{j}| < |v_{j+1}|$ for every integer $\frac{n}{2} < j < n$. Let $\delta$ be defined by $v_j = G_n(j + 1 - \delta)$, we can compute
  \begin{align*}
    \frac{|G_n(j+1 - \delta)|}{|G_n(j+1 + \delta)|} &= \frac{(j-\delta)(j-1-\delta)\cdots(1-\delta)\cdot \delta(1+\delta)\cdots(n-j-1+\delta)}{(j+\delta)(j-1+\delta)\cdots \delta \cdot (1-\delta)(2-\delta)\cdots(n-j-1-\delta)}\\
    &=\frac{(n-j-\delta)(n-j+1-\delta)\cdots(j-\delta)}{(n-j+\delta)(n-j+1+\delta)\cdots(j+\delta)} < 1.
  \end{align*} This implies $|v_{j}| = |G_n(j+1 - \delta)| < |G_n(j+1 + \delta)| \leq |v_{j+1}|$ as wanted.
\end{proof}

\end{document}
